\documentclass[11pt,a4paper]{article}
\usepackage[utf8]{inputenc}
\usepackage[T1]{fontenc}
\usepackage[english]{babel}
\usepackage{enumitem}
\usepackage[hidelinks]{hyperref}

\title{How to Work With This Thesis Template}
\author{}
\date{}

\begin{document}
\maketitle

\href{https://www.overleaf.com/read/fxhpwgkznxdx#ca782c}{Link to this project}

\section{Overview}
This project is a thesis template built around \texttt{main.tex}. The main file pulls in the cover pages, the front matter, the chapter list, the bibliography, and the Hebrew table of contents. The files you will edit most often are:
\begin{itemize}[leftmargin=*]
    \item \texttt{thesis\_metadata.tex} for title, author, advisors, and date.
    \item \texttt{thesis\_chapters.tex} for chapter order and appendix order.
    \item \texttt{content/.../*.tex} for the actual chapter text.
\end{itemize}

\section{Fill In Metadata}
Open \texttt{thesis\_metadata.tex} and replace the placeholder values with your thesis details.

\begin{itemize}[leftmargin=*]
    \item Update the English and Hebrew thesis titles.
    \item Update the faculty name if needed.
    \item Update the thesis program name.
    \item Replace the author's name with your own name in both languages.
    \item Replace the advisor fields. If you have more than one advisor, separate the English names with \texttt{\string\and}.
    \item Set the submission month and year.
\end{itemize}

These values are reused across the cover and abstract pages, so changing them in one place updates the entire document.

\section{Fill In Chapter Names and Order}
Open \texttt{thesis\_chapters.tex} to control the thesis structure.

Each chapter is declared with:
\begin{verbatim}
\thesischapter{English Title}{chapter:label}
              {Hebrew Title}{content/path/file}
\end{verbatim}

Each appendix is declared with:
\begin{verbatim}
\thesisappendix{English Title}{appendix:label}
               {Hebrew Title}{content/path/file}
\end{verbatim}

\begin{itemize}[leftmargin=*]
    \item Keep the chapter order in the order you want it to appear in the thesis.
    \item Add new chapters by creating a new file under \texttt{content/} and adding a matching \texttt{\string\thesischapter} line.
    \item Remove chapters you do not need by deleting or commenting out their lines.
    \item If you do not need appendices, remove the \texttt{appendices} block.
\end{itemize}

The provided chapter files are already wired into \texttt{main.tex}, so the thesis will include whatever you list here.

\section{Add Content}
\begin{itemize}[leftmargin=*]
    \item Update the English and Hebrew abstract files in \path{Main_pages/abstract.tex} and \path{Main_pages/abstract_heb.tex}.
    \item Update the acknowledgements file in \path{Main_pages/acknowledgements.tex}.
    \item Add content to the chapter files under \texttt{content/}.
    \item Add appendices in \path{content/appendices/} if needed.
\end{itemize}

A few practical tips:
\begin{itemize}[leftmargin=*]
    \item Use \texttt{\string\label} and \texttt{\string\ref} for cross-references.
    \item Add citations with \texttt{\string\citep} or \texttt{\string\citet} and keep references in the bibliography file.
    \item Keep figures in the chapter folders or in \path{content/images/} and reference them with \texttt{\string\includegraphics}.
    \item Compile \texttt{main.tex} after major edits to keep the table of contents, bibliography, and Hebrew TOC in sync.
\end{itemize}

\section{Build Flow}
The usual workflow is:
\begin{enumerate}[leftmargin=*]
    \item Edit \texttt{thesis\_metadata.tex}.
    \item Edit \texttt{thesis\_chapters.tex}.
    \item Fill in the chapter files under \texttt{content/}.
    \item Compile \texttt{main.tex} until the references and tables of contents settle.
\end{enumerate}

If you see old table-of-contents or Hebrew TOC entries after structural changes, delete the auxiliary build files and compile again.

\end{document}
